\documentclass[%
 reprint,
 amsmath,amssymb,
 aps,
]{revtex4-1}

\usepackage{graphicx}% Include figure files
\usepackage{dcolumn}% Align table columns on decimal point
\usepackage{bm}% bold math
\usepackage[dvipdfm,colorlinks=true,linkcolor=blue,filecolor=blue, urlcolor=blue]{hyperref}
\usepackage{braket}
\usepackage{caption} % for captionsetup[figure]
\usepackage{threeparttable} % for table caption modification
\usepackage{breqn} % dmath: automatic equation line breaking
\numberwithin{equation}{section}% (p.37)

% \sloppy sets all text to loose spacing, allowing overflow and uneven word spacing;
% not recommended for academic papers requiring strict typesetting.
\tolerance=1000
\emergencystretch=1em
\hyphenpenalty=3000
\exhyphenpenalty=3000
\flushbottom
\usepackage[final]{microtype}


\begin{document}

    \captionsetup[figure]{justification=raggedright,
        labelfont={bf},name={Fig.},labelsep=period} % raggedright for left-aligned figure captions
    \captionsetup[table]{
        labelfont={bf},name={Table.},labelsep=period}


\title{Spin-Induced Inertial Resistance in Electrons:\\A Gyroscopic Interpretation Based on General Relativity}

\author{Satoshi Hanamura}
\email{hana.tensor@gmail.com}
\date{April 20, 2025}

\begin{abstract}
This study provides substantial support for the theory of Zitterbewegung within the electron and its association with gyroscopic effects. While Zitterbewegung has traditionally been interpreted as an oscillatory term in the Dirac equation, we have geometrically reconstructed this phenomenon through our model of spatially separated energy kernels exchanging thermal potential energy. Applying special relativistic principles and Lorentz transformation to this model yields an electron Zitterbewegung velocity of 0.040472$c$. Furthermore, by algebraically incorporating general relativistic effects through geodetic precession, we refine this velocity to 0.040374$c$.

A key contribution of this study is the introduction of a remarkably simple and dimensionally consistent equation,
\[
\gamma = 1 + a,
\]
which relates the dimensionless anomalous magnetic moment of the electron to the Lorentz factor from special relativity. This compact formulation captures the essence of both quantum correction and relativistic kinematics, establishing a direct correspondence between two foundational pillars of modern physics.

By interpreting spin as a deterministic oscillatory motion within the electron—analogous to a relativistic harmonic oscillator—we derive a geometric model in which internal angular momentum gives rise to gyroscopic resistance. Just as classical gyroscopes resist directional changes due to their spin, we propose that this internal structure underpins the electron’s inertial mass. The model offers a unified perspective in which a single electron, through its intrinsic spin dynamics, exhibits resistance to acceleration consistent with classical inertia.
\end{abstract}

\maketitle

% ========================================
% Section 1: Introduction
% ========================================
\section{Introduction}

In conventional quantum theory, the electron has traditionally been treated as a dimensionless point particle. This mathematical abstraction has served physics well in many contexts, but it also leads to well-known theoretical difficulties, such as infinite self-energy and the challenges of quantum field renormalization.

To facilitate a paradigm shift in our understanding, let us begin with an intuitive analogy:

\begin{quote}
\textit{Imagine the electron not as a featureless point, but as a sophisticated entity with internal structure—like a miniature clock with intricate gears. These ``gears'' represent the internal oscillatory motion (Zitterbewegung) that gives the electron its intrinsic properties. This internal mechanism not only measures time locally but also interacts with surrounding spacetime. Just as a mechanical watch both measures and is affected by its environment, the electron's internal oscillations both define and respond to spacetime properties. This conceptual reframing might provide the key to reconciling quantum mechanics with general relativity—two theoretical frameworks that have remained stubbornly incompatible for nearly a century.}
\end{quote}

This analogy, while simplified, captures the essence of our theoretical proposal: that elementary particles possess internal spatiotemporal structure which defines their properties and relationships with the broader universe. Instead of treating the electron as a mathematical point responding to external spacetime, we propose that the electron's internal oscillations are themselves manifestations of localized spacetime structure.

Building on our work~\cite{hanamura2023}, this study develops a deterministic alternative to the standard model, describing electrons as closed systems with conserved energy, independent of quantum field theory and the Heisenberg uncertainty principle. By avoiding probabilistic collapse, our model aligns with deterministic interpretations~\cite{bohm1952}. Our 0-Sphere model reinterprets Zitterbewegung, traditionally a mathematical artifact of the Dirac equation~\cite{schrodinger1930}, as a deterministic oscillation between two energy kernels. By integrating special and general relativity, we predict a Zitterbewegung velocity of $0.040374c$, refined via geodetic precession, and propose that electron spin, analogous to a classical gyroscope's resistance, generates inertial mass.

In classical mechanics, gyroscopes exhibit resistance to changes in their spin axis due to the conservation of angular momentum—a well-established phenomenon known as the gyroscopic effect. While quantum spin does not correspond to literal rotational motion, it shares mathematical properties with classical angular momentum. In this report, we propose a novel interpretation of electron spin as a microscopic analog of a gyroscope, where resistance to directional change arises from intrinsic spin dynamics modeled via relativistic precession.

This work builds on our previous theoretical framework, in which the electron is modeled as a harmonic oscillator composed of two energy kernels that exchange thermal potential energy via a sinusoidal mechanism. Analogous to a miniature internal clock, this model reinterprets the Dirac equation by viewing its positive and negative energy solutions as coexisting components of a single electron. The resulting internal oscillation (Zitterbewegung) embodies a deterministic structure that unifies quantum and relativistic phenomena.

% ========================================
% Section 2: Theoretical Framework for Spin-Induced Inertia
% ========================================
\section{Theoretical Framework for Spin-Induced Inertia}

% ----------------------------------------
% Subsection 2.1: Mathematical Formulation of the 0-Sphere Model
% ----------------------------------------
\subsection{Mathematical Formulation of the 0-Sphere Model}
\label{ref:MathematicalFormulation}

The name ``0-Sphere model'' derives from the geometric concept of a 0-sphere, which represents two discrete points in space. These two points correspond to Kernel $A$ and Kernel $B$ in our model.

This duality between Kernel $A$ and Kernel $B$ is conceptually supported by our reinterpretation of the Dirac equation: we assume that a single electron inherently incorporates both the positive and negative energy solutions. Rather than treating these solutions as abstract mathematical constructs or virtual states, we propose that they correspond to real, coexisting internal components---Kernel $A$ and Kernel $B$---within the electron. Specifically, the kernel that emits Thermal Potential Energy (TPE) is associated with the negative energy solution, while the kernel that absorbs TPE corresponds to the positive energy solution. This thermodynamic interpretation provides a physical mechanism for embedding an internal oscillatory clock within the electron, and explains the emergence of Zitterbewegung as a manifestation of this internal duality.

In the 0-Sphere electron model, the electron's mass can be converted to thermal energy, which we designate as TPE. At position $+a$, we place TPE $A$, referred to as Kernel $A$, representing the electron's mass. This TPE transforms from mass energy to radiation energy as it moves toward position $-a$. This process can be represented as:

\begin{equation}
T_{\mathrm{kernel}A} \rightarrow \gamma^*_{\mathrm{K.E.}} \rightarrow T_{\mathrm{kernel}B},
\label{eq:electron_transfer}
\end{equation}
where $T_{\mathrm{kernel}A}$ and $T_{\mathrm{kernel}B}$ denote the thermal potential energies associated with Kernel $A$ and Kernel $B$, respectively. This radiation energy is converted into kinetic energy, which acts as the driving force for the spatial motion of the photon sphere that connects Kernel $A$ and Kernel $B$. This dynamic structure is referred to as the ``$\gamma^*_{\mathrm{K.E.}}$ : photon sphere''. The temperature gradient generated by Kernels $A$ and $B$ produces a clean sinusoidal wave~\cite{hanamura2018}:

\begin{equation}
\mathrm{grad} \ (T_{\mathrm{kernel}B} - T_{\mathrm{kernel}A}) \propto \sin (\omega t) .
\label{e3_simple_harmonic_motion}
\end{equation}

The energy transportation pathway between Kernel $A$ and Kernel $B$ via thermal radiation follows a geodesic trajectory. This establishes a model that coherently connects quantum oscillations with general relativity, as radiation energy naturally follows geodesic paths. Notably, this radiative energy transfer is further constrained by Snell's law, which governs the geometric trajectory of the photon sphere, as explored in detail in our work \cite{hanamura2024}. This perspective offers a unified framework for understanding both the quantum oscillations within the electron and the geometric structure of spacetime, potentially bridging the conceptual gap between quantum mechanics and general relativity.

However, according to Stefan-Boltzmann's law, the relationship between radiation energy $I$ and temperature $T$ is:

\begin{equation}
I = \sigma T^4 .
\end{equation}

We model Kernels \( A \) and \( B \) as simple harmonic oscillators modeled using cosine and sine functions, respectively. Their corresponding radiation energies are each raised to the fourth power in accordance with the Stefan–Boltzmann law. Additionally, by including the squared sine term corresponding to the kinetic energy governed by the photon sphere, the internal structure of the electron satisfies a clean energy conservation law. Specifically, the sum of the three terms within the parentheses on the right-hand side of the equation below is identically equal to unity.

When incorporating the kinetic energy associated with the photon sphere, the internal structure of a single electron can oscillate while conserving its total energy.  
Let \( E_0 \) denote the rest energy of the electron, i.e., \( E_0 = mc^2 \). Under this assumption, the following identity holds:

\begin{equation}
E_{0}(t) = E_{0} \left( \cos^{4} \left( \frac{\omega t}{2} \right) + \sin^{4} \left( \frac{\omega t}{2} \right) + \frac{1}{2} \sin^{2} (\omega t) \right).
\label{eq:total_hamiltonian}
\end{equation}

\vspace{0.3cm}

Our model posits that, due to the time phase, there exists a specific phase in which the total energy of a single-electron system is entirely localized within Kernel $A$. At this phase, both the thermal potential energy associated with Kernel $B$ and the kinetic energy of the photon sphere become zero. Specifically, when all of the electron's rest mass becomes concentrated in Kernel $A$, the cosine term in Eq.~\ref{eq:total_hamiltonian} equals 1, while the second and third terms (the sine terms) on the right-hand side become zero. Conversely, there exists another phase in which the total rest mass energy is entirely localized within Kernel $B$.

At this complementary phase, the first cosine term and the third sine-squared term on the right-hand side of the equation become zero, indicating that both the thermal potential energy associated with Kernel $A$ and the kinetic energy of the photon sphere become zero. To represent this behavior, we adopted a combination of sine and cosine functions, modeling the electron as an oscillator. Furthermore, to incorporate the double-valued nature of spinors, the thermal potential energy is expressed along both the real and imaginary axes. This mathematical structure is expected to offer extensibility in the forthcoming theoretical developments.


% ----------------------------------------
% Figure 1: f-TotalHamiltonian.png -- Energy conservation in the 0-Sphere model
% ----------------------------------------
\begin{figure*}[t!]
    \centering 
    \includegraphics[width=\textwidth]{f-TotalHamiltonian.png}
    \caption{Energy conservation in the 0-Sphere electron model. The graph illustrates the time evolution of energy components: thermal potential energy (TPE) in Kernel $A$ (blue dashed line, $\cos^4(\phi/2)$), TPE in Kernel $B$ (orange dotted line, $\sin^4(\phi/2)$), and kinetic energy of the photon sphere (green dotted line, $(1/2)\sin^2(\phi)$). At $\phi=0$, all rest energy is in Kernel $A$; at $\phi=\pi$, all rest energy transfers to Kernel $B$; and at $\phi=2\pi$, the cycle completes back to Kernel $A$. The red solid line ($H(\phi)=1$) demonstrates that total energy remains constant throughout the oscillation cycle. This visualization represents the energy transfer process $T_{\mathrm{kernel A}} \rightarrow \gamma^*_{\mathrm{K.E.}} \rightarrow T_{\mathrm{kernel B}}$ where the spatially separated kernels form a 0-sphere structure.}
    \label{f-TotalHamiltonian}
\end{figure*}


% ----------------------------------------
% Subsection 2.2: The Seesaw and Basketball Analogy
% ----------------------------------------
\subsection{The Seesaw and Basketball Analogy for Simple Harmonic Motion}

% ----------------------------------------
% Subsubsection 2.2.1: Conceptual Visualization of the 0-Sphere Model
% ----------------------------------------
\subsubsection{Conceptual Visualization of the 0-Sphere Model}
To better conceptualize the 0-Sphere electron model~\cite{hanamura2020}, let us consider a seesaw at a playground, tilted to one side. Imagine placing a basketball at the higher end of the seesaw. If you stabilize the seesaw with your hand, the basketball would roll down toward the lower end and eventually fall off, continuing to roll on the ground. However, if you manipulate the seesaw at appropriate intervals, you could cause the basketball to travel back and forth in a controlled manner. This oscillatory motion forms the conceptual basis for simple harmonic oscillation.

Now, imagine that you and a friend sit on opposite ends of the seesaw. There exists a physical distance between you and your friend—this represents discreteness. Meanwhile, the basketball rolls continuously from your position to your friend's position. The 0-Sphere model embodies this coexistence of continuous and discrete elements.

To illustrate this state, we can compare it to Alice and Bob sitting on opposite sides of a seesaw, passing a basketball back and forth. Here, Alice represents Kernel $A$, and Bob represents Kernel $B$. When the $\cos^4(\omega t/2)$ term equals 1, the basketball is entirely in Alice's hands, indicating that the total energy is localized in Kernel $A$. Conversely, when the $\sin^4(\omega t/2)$ term equals 1, the basketball has reached Bob's hands, meaning the energy is localized in Kernel $B$. In both of these extreme phases, the basketball is stationary, and the kinetic energy of the basketball---comparable to the kinetic energy of the photon sphere---is zero.

However, during the intermediate phases, the basketball is in motion and is in neither Alice's nor Bob's hands. These moments correspond to the activation of the $\frac{1}{2} \sin^2(\omega t)$ term in the energy identity, representing the kinetic energy of the photon sphere. In this way, the time evolution of the system not only describes the periodic exchange of thermal potential energy between Kernel $A$ and Kernel $B$, but also captures the transient manifestation of kinetic energy that mediates between them.

% ----------------------------------------
% Subsubsection 2.2.2: Mathematical Representation of Energy Conservation
% ----------------------------------------
\subsubsection{Mathematical Representation of Energy Conservation}
The use of fourth-power terms in Eq.~\eqref{eq:total_hamiltonian} is motivated by the Stefan–Boltzmann law, where radiative energy is proportional to the fourth power of temperature. In this model, the sinusoidal temperature distributions of Kernel $A$ and Kernel $B$ are expressed as $\cos^4(\omega t/2)$ and $\sin^4(\omega t/2)$, whose fourth powers represent the respective energy contributions. The inclusion of the kinetic energy term, $\frac{1}{2}\sin^2(\omega t)$, ensures total energy conservation throughout the oscillatory cycle. The mathematical identity:
\[
\cos^4\left(\frac{\omega t}{2}\right) + \sin^4\left(\frac{\omega t}{2}\right) + \frac{1}{2}\sin^2(\omega t) = 1,
\]
demonstrates that the sum of these components remains constant, validating the internal energy balance of the electron structure. This energy conservation and the periodic exchange between different energy forms is visualized in Figure~\ref{f-TotalHamiltonian}.

In this model, a closed algebraic equation is used to predict the energy distribution of an electron without relying on perturbation theory, offering a deterministic view of quantum phenomena. To analytically represent this behavior, the model describes the total energy using a sinusoidal function, with kinetic energy and thermal potential energy exchanging periodically over time.

A remarkable feature of this equation is that while the photon sphere completes one full cycle with a 360° phase change, the internal structures—referred to here as kernels $A$ and $B$—require a 720° phase change to complete their respective cycles. This provides a mathematical explanation for the difference in spin: photons exhibit spin 1 and complete a full period with 360°, whereas electrons possess spin 1/2 and require 720° for a full cycle.


% ----------------------------------------
% Subsection 2.3: Spin as a Relativistic Harmonic Oscillator
% ----------------------------------------
\subsection{Spin as a Relativistic Harmonic Oscillator}

\textit{\small This section omits the mathematical derivations. For any unclear points regarding derivation methods or procedures, please refer to our previous papers as appropriate.}

\vspace{0.5cm}

Electron spin has long been interpreted as an intrinsic form of angular momentum. In our recent theory \cite{hanamura2023}, we redefine spin as arising from sinusoidal motion between two energy kernels, yielding angular velocity governed by the Thomas precession as first noted by Thomas \cite{thomas1926} and later developed in full \cite{thomas1927}:
\begin{equation}
\boldsymbol{\Omega} = \frac{1}{2c^2} (\mathbf{a} \times \mathbf{v}).
\end{equation}

When both acceleration $\mathbf{a} = -\sin\omega t$ and velocity $\mathbf{v} = \cos\omega t$ are sinusoidal, their cross product results in an oscillating angular velocity with a double frequency:
\begin{equation}
\boldsymbol{\Omega}(t) = \frac{1}{2c^2} \cdot \left(-\frac{1}{2}\sin2\omega t\right).
\end{equation}

Unlike the standard model's field-theoretic spin operators~\cite{weinberg1995}, our model derives spin from deterministic oscillations, offering a geometric alternative. This intrinsic precession leads to ``a repeating up/down spin state'' and underpins the anomalous magnetic moment via Lorentz contraction. The doubling of frequency provides a classical basis for the quantization of spin angular momentum to half the Planck constant.

To intuitively illustrate this internal structure, we introduce the analogy of a seesaw with a basketball rolling between two ends. When the ball is in Alice's hands (Kernel $A$), all energy is localized there; when it reaches Bob (Kernel $B$), the energy is transferred. Between these extremes, the ball is in motion—analogous to the kinetic energy of the photon sphere that mediates energy exchange. This discrete-yet-continuous motion encapsulates the dynamics of the 0-Sphere model~\cite{hanamura2024}.

As shown in Figure~\ref{f-TotalHamiltonian}, the time evolution of energy components within the 0-Sphere model follows a precise conservation law. The rest energy oscillates between two thermal potential energy kernels while the kinetic energy of the connecting photon sphere mediates the transition, forming a stable, periodic structure.

We argue that the anomalous magnetic moment of the electron can be understood through rotational Lorentz contraction. In a rotating coordinate system, Einstein noted that ``the ratio of circumference to diameter deviates from $\pi$~\cite{einstein1912}.'' When applied to electron spin, this suggests that a rotation affected by Lorentz contraction becomes shorter, with the difference interpreted as the anomalous magnetic moment:
\begin{equation}
\frac{L}{L_0} = \frac{1}{1 + \frac{1}{\sqrt{2}}a_e^{\text{exp}}},
\label{eq:zitter_velocity}
\end{equation}
where $a_e^{\text{exp}}$ is the experimentally measured anomalous magnetic moment~\cite{fan2023}.

The formula above establishes a profound connection between quantum corrections (represented by the anomalous magnetic moment $a$) and relativistic kinematics (represented by the Lorentz factor). It calculates the average oscillation velocity of Zitterbewegung, which explains the inclusion of the root-mean-square (RMS) term. This use of the RMS value is motivated by the physical analogy that the force driving the Zitterbewegung is similar to the driving force in an AC harmonic oscillator, where the effective value is naturally represented by the RMS of the oscillation. For expressing the maximum velocity of Zitterbewegung oscillation, denoted as $v_{\mathrm{ZBmax}}$, which represents the maximum internal oscillation velocity of the electron within its own rest frame. Using the well-known relationship $\gamma = 1/\sqrt{1 - \beta^2} = 1/\sqrt{1 - v_{\mathrm{ZBmax}}^2/c^2}$, we arrive at:

\begin{equation}
\gamma = 1 + a .
\label{eq:gamma_anomalous}
\end{equation}

\vspace{0.3cm}


This simple relationship serves as ``a bridge connecting the quantum mechanical micro-oscillations of the electron with special relativity'', providing a unified framework for understanding these previously separate physical domains.

This framework, based solely on special relativistic calculations, initially yielded the remarkable result that a stationary free electron's average trembling motion velocity within Compton wavelengths is approximately $0.040472c$ or about 12,142 km/s. However, our research has developed a model that applies the concept of geodetic precession from general relativity to quantum oscillations, leading to a refined prediction.

% ----------------------------------------
% Subsection 2.4: Geodetic Precession and General Relativistic Corrections
% ----------------------------------------
\subsection{Geodetic Precession and General Relativistic Corrections}

Just as a classical gyroscope resists reorientation, the electron's spin axis—interpreted here as an outcome of continuous precessional motion—should exhibit resistance to external disturbances. This is especially relevant under acceleration, where relativistic effects become significant.

According to general relativity, the geodetic precession experienced by a mass in orbit around a central mass is given by~\cite{hartle2021}:
\begin{equation}
\Delta\varphi_{\text{geodetic}} = 2\pi \left[1 - \left(1 - \frac{3M}{R} \right)^{1/2} \right].
\label{eq:geodetic_precession}
\end{equation}

Applied to the electron modeled as a harmonic oscillator between discrete kernels, this implies that even microscopic spin systems experience a curved-space-induced precession. By incorporating geodetic effects, we can modify our equation for Lorentz contraction:

\vspace{0.1cm}

\begin{equation}
\frac{L}{L_0} = \frac{1}{1 + \frac{1}{\sqrt{2}}a_e^{\text{exp}} - \frac{\Delta\varphi_{\text{geodetic}}}{2\pi}}.
\label{eq:core_equation}
\end{equation}

\vspace{0.3cm}

When we incorporate these general relativistic corrections through geodetic precession, our prediction for the electron's Zitterbewegung velocity is refined from the purely special relativistic value of $v_{e,\text{SR}} = 0.040472c$ to $v_{e,\text{SR+GR}} = 0.040374c$. This small but significant difference demonstrates the subtle interplay between quantum mechanics and general relativity at microscopic scales.

This approach allows us to estimate the electron's radius by measuring its trembling motion velocity, potentially bridging quantum mechanics and general relativity. Our calculations suggest that if the electron's velocity is approximately $0.04c$, its radius would be between $10^{-25}$ and $10^{-26}$ meters~\cite{hanamura2025}.


% ----------------------------------------
% Subsection 2.5: Integration of Proper Time into the Particle Model
% ----------------------------------------
\subsection{Integration of Proper Time into the Particle Model}

The energy distribution equation (Eq.~\ref{eq:total_hamiltonian}) introduced in Section ~\ref{ref:MathematicalFormulation} has an important property: it algebraically permits the assignment of a unique time parameter $t$ to each individual electron. This suggests the possibility of incorporating the general relativistic concept of ``proper time directly into the particle model''. Furthermore, the formulation can be extended to account for gravitational effects, where oscillation frequencies decrease under gravitational influence—a phenomenon observable in real elementary particles. Consequently, the time parameter $t$ in this equation can be replaced with proper time $\tau$, providing a more comprehensive representation that accommodates gravitational field effects.

This conceptual framework enables quantitative predictions regarding electron behavior under varying gravitational field conditions. For instance, if measurements of changes in electron Zitterbewegung within strong gravitational fields could be obtained, they would provide a direct means of testing this theoretical model. Such experiments would offer compelling evidence for the integration of proper time into the particle's intrinsic structure and validate the proposed unification of quantum mechanics and general relativity.

% ----------------------------------------
% Subsection 2.6: Electron-Positron Annihilation
% ----------------------------------------
\subsection{Electron-Positron Annihilation in the Spacetime Oscillator Framework}
\label{Electron-Positron}

If the electron represents a stable state in this framework due to its critical radius properties under normal conditions, this approach offers insights into why nature exhibits three generations of leptons with their characteristic mass differences. The substantial gaps observed in lepton mass hierarchy may be understood through the corresponding gaps in critical radii at which internal Zitterbewegung dynamics become unsustainable. According to our view, tau particles and muons are subject to mechanisms that decrease their critical radius, causing them to decay into particles of lower hierarchies. This decay occurs when the particle's radius reaches the critical radius, at which point the Zitterbewegung oscillation velocity becomes zero, triggering decay into lower-level particles. However, the electron typically does not reach its critical radius and thus maintains a stable state. 

When applying our model to observed phenomena, particularly the electron-positron annihilation, a more profound physical mechanism emerges. Upon collision, the interaction between electrons and positrons likely reduces their Zitterbewegung to zero, effectively nullifying the oscillation between kernels $A$ and $B$ in both particles. This cessation of internal oscillation would lead to ``the release of the photon sphere structure''---which mediates energy transfer between the kernels in our model—manifesting as the observable photons produced during annihilation. This interpretation transforms our understanding of annihilation from an abstract quantum field theoretical process to a concrete geometric event involving specific internal structures. Furthermore, this model might enable new predictions regarding the characteristics of emitted photons during annihilation, such as their polarization states or angular distribution patterns, potentially offering experimental verification pathways for the theory.

Rather than merely explaining the exact number of generations, this model provides a geometric mechanism that connects the discrete mass transitions in lepton decays to quantized critical radii determined by the underlying spacetime geometry. The unified framework capable of addressing both lepton decay hierarchies and electron-positron annihilation demonstrates the conceptual coherence and explanatory power of viewing particles as localized spacetime oscillators rather than abstract point entities existing within spacetime.


% ========================================
% Section 3: Gyroscopic Analogy and the Origin of Inertial Mass
% ========================================
\section{Gyroscopic Analogy and the Origin of Inertial Mass}

% ----------------------------------------
% Subsection 3.1: Theoretical Framework for Gyroscopic Inertia
% ----------------------------------------
\subsection{Theoretical Framework for Gyroscopic Inertia}

% ----------------------------------------
% Subsubsection 3.1.1: Geometric Equations Linking ZB and Relativity
% ----------------------------------------
\subsubsection{Geometric Equations Linking Zitterbewegung and Relativity}

Building on the premise that the electron possesses an internal structure rather than being a point-like entity, we propose that its properties evolve with time in a periodic fashion. This temporal evolution is represented by Eq.~\eqref{eq:electron_transfer}. The thermal potential energy and kinetic energy, both of which fluctuate according to the temporal phase $E_{0}(t)$ of the rest energy, satisfy the law of energy conservation within the system of a single electron. This is expressed by Eq.~\eqref{eq:total_hamiltonian}.

The quartic sine and cosine terms in Eq.~\eqref{eq:total_hamiltonian} represent potential energy. When these terms are differentiated with respect to time and their difference is taken, the resulting expression reduces to a simple harmonic oscillator described by $\sin(\omega t)$. This implies that the force $\bm{\mathrm{F}}$ responsible for activating the photon sphere behaves as a harmonic oscillator, as derived in Eq.~\eqref{e3_simple_harmonic_motion}. For detailed derivations of the radiation gradient that imparts kinetic energy to the photon sphere, the reader is referred to our previous works~\cite{hanamura2018, hanamura2024}, where the TPE difference between the two energy kernels gives rise to the oscillatory driving force.

The considerations and derivations presented in the previous sections were developed in our earlier publications. The novel contribution of this study lies in the proposed mechanism whereby the electron acquires inertial mass through internal dynamical resistance—a mechanism absent from conventional field-theoretic or relativistic frameworks.

As discussed in the previous section, Eq.~(\ref{eq:zitter_velocity}) predicts the Zitterbewegung velocity of the electron. The factor of $\sqrt{2}$ appears due to the averaging over the sinusoidal motion. For the maximum velocity in simple harmonic oscillation, the corresponding expression becomes Eq.~(\ref{eq:gamma_anomalous}). This equation is remarkably elegant, as it suggests the equivalence between the Lorentz transformation in special relativity and the anomalous magnetic moment.

Equation~(\ref{eq:geodetic_precession}) describes geodetic precession and is a standard result found in general relativity textbooks. Equation~(\ref{eq:core_equation}) forms the core of our theoretical framework: it places Lorentz contraction on the left-hand side and incorporates Eq.~(\ref{eq:zitter_velocity})—which expresses Lorentz contraction equivalently—on the right-hand side, while also correcting it with a geodetic precession term from general relativity. In this sense, Eqs.~(\ref{eq:zitter_velocity}, \ref{eq:gamma_anomalous}), and (\ref{eq:core_equation}) are newly proposed geometric equations that we have developed to bridge relativity and quantum mechanics.


% ----------------------------------------
% Subsubsection 3.1.2: Gyroscopic Resistance as the Origin of Inertial Mass
% ----------------------------------------
\subsubsection{Gyroscopic Resistance as the Origin of Inertial Mass}

Classical gyroscopes resist changes to their orientation due to the conservation of angular momentum. When a torque is applied perpendicular to the spin axis, the gyroscope responds with precession—a motion perpendicular to both the torque and the spin axis. This resistance to directional change gives gyroscopes their characteristic ``perceived resistance'' when manipulated. We suggest that the familiar resistance experienced when manipulating a gyroscope---arising from its internal angular momentum---is fundamentally analogous to the origin of inertial mass at the quantum level. Similarly, if electrons possess intrinsic spin angular momentum as described in our model, they would exhibit an analogous resistance to acceleration—a property we interpret as inertial mass.

The gyroscopic effect is a physical phenomenon where a rotating object tends to maintain its axis of rotation. In conventional physics, this effect is based on the conservation of angular momentum generated by rotational motion. However, in this study's 0-Sphere model, this concept is fundamentally reconsidered. The electron's spin in this model is reinterpreted as a back-and-forth oscillation between Kernels $A$ and $B$, representing a linear vibration rather than the traditional circular motion. This revision of the conventional concept that spin angular momentum results from circular motion is a key innovation of our model. In essence, each electron has an embedded mechanism that tries to maintain the axis formed by Kernels $A$ and $B$, and this internal structure generates resistance to external acceleration—what we perceive as inertial mass.

A critical insight emerges when examining the energetic pathway between Kernels $A$ and $B$ during non-zero kinetic energy phases. The directional radiation and absorption of thermal energy between these kernels establishes a virtual axis within the electron---not a physical rod, but a preferred direction of energy transport. This axis is fundamental to understanding the gyroscopic behavior of the electron.

An illustrative analogy can be drawn from celestial mechanics: in the 19th century, astronomers observed that Uranus’ orbit exhibited slight but persistent anomalies---wobbles that could not be explained by Newtonian mechanics alone. These irregularities were ultimately attributed to the gravitational influence of an unseen planet---later identified as Neptune---exerting directional perturbations on Uranus’ orbit.

Similarly, in our model, the electron’s internal energy axis resists directional changes induced by external forces. When an external force attempts to change the electron's orientation, this internal energy transport pathway must also reorient. However, due to energy conservation principles, this reconfiguration encounters resistance. The system actively opposes changes to its established energy exchange trajectory, manifesting as an inertial resistance to acceleration. This resistance mechanism provides a geometric explanation for inertial mass: rather than being an intrinsic or externally bestowed property, mass emerges as the resistance to altering established internal energy transport pathways. This perspective unifies the concepts of spin, inertia, and mass through the electron's internal structural dynamics, offering a deterministic alternative to field-theoretic interpretations.

% ----------------------------------------
% Subsection 3.2: Connections to Fundamental Physics and Broader Implications
% ----------------------------------------
\subsection{Connections to Fundamental Physics and Broader Implications}

This interpretation aligns with Penrose's argument that spin is not merely an internal degree of freedom but plays a fundamental role in shaping space-time structure. In his combinatorial approach, Penrose suggests that the angular momentum of particles, particularly spin-$\frac{1}{2}$ systems, serves as a foundational element for constructing space-time itself. In our model, the intrinsic Zitterbewegung-induced spin of the electron similarly generates a local geometric framework—analogous to Penrose’s ``twistor space''—which gives rise to inertial properties through its resistance to directional changes. This conceptual bridge reinforces the view that spin and geometry are deeply interconnected~\cite{penrose1971}.

To further elaborate on this gyroscopic mechanism of inertial mass, it is crucial to recognize that the resistance exhibited by a classical gyroscope requires no external energy input. This resistance emerges purely from the geometric configuration and motion state of the internal structure. When we attempt to change the direction of a gyroscope's angular momentum vector, we experience a resistance that is proportional to both the magnitude of the angular momentum and the rate at which we try to change its direction.

By analogy, in our electron model, the Zitterbewegung oscillation establishes an intrinsic angular momentum that resists changes in its motion state. The faster the internal Zitterbewegung velocity, the greater the resistance to acceleration—which manifests as a larger inertial mass. This suggests a potential explanation for the mass hierarchy observed among leptons, where heavier leptons like muons and tau particles might possess correspondingly different intrinsic Zitterbewegung dynamics.

As noted in subsection~\ref{Electron-Positron}, we propose that when an electron and positron annihilate, their respective kernels cancel each other out. This cancellation halts the radiation and absorption cycle between Kernels $A$ and $B$, effectively eliminating the Zitterbewegung. As a result, only the photon sphere remains, manifesting as gamma rays emitted in the aftermath of the annihilation. Extending this concept further, the reason photons possess no mass can be derived from the absence of internal Zitterbewegung. This perspective offers a geometric explanation for both the presence of mass in electrons and its absence in photons, unifying our understanding of these fundamental particles through their internal structural dynamics. This geometric interpretation might offer a perspective where inertial mass emerges naturally from internal structural dynamics rather than being an inherent or externally imposed property. 

This view has interesting connections to Mach's principle~\cite{mach1919}, which proposes that inertia does not exist in isolation but arises from the relationship between a body and the rest of the universe. While Mach considered these relationships in terms of distant masses, our model suggests an alternative where the relationship exists within the particle's own internal geometry. This distinction is critical: whereas Mach's principle attributes inertia to global interactions with the distant cosmos, our model posits a self---contained, local mechanism. Although both frameworks challenge the notion of inertia as an intrinsic property, they do so from fundamentally different relational perspectives---external versus internal.

Therefore, while our model resonates with Mach’s broader philosophy of relational origins of inertia, it deviates from its cosmological context by grounding inertia in localized internal motion rather than distant mass distributions.
The Zitterbewegung oscillation creates an internal reference system—a form of ``local universe''—against which inertial resistance manifests. Though distinct from Mach's original conception—as reinterpreted in a modern context by Barbour~\cite{barbour2001}—both approaches share the fundamental insight that inertia may not be an intrinsic property but emerges from structural relationships.

Similarly, this framework provides a geometric foundation that could potentially contribute to understanding the observed equivalence between inertial and gravitational mass. Since both the electron's internal Zitterbewegung and its gravitational interactions involve curvature in spacetime geometry, albeit at different scales, a deeper connection may exist. However, a complete reconciliation with Einstein's equivalence principle would require further theoretical development linking these internal dynamics to the particle's gravitational behavior.

Departing from the traditional point-particle view, this model envisions the electron as possessing an independent oscillatory structure—an internal mechanism analogous to a cosmic gyroscope. Such a framework provides a new particle image in which each electron functions as an autonomous, structured system. The central conclusion derived from the equations presented in this paper is that the electron, as described here, does not require energy fluctuations to manifest inertial resistance; rather, it arises naturally from the proposed internal dynamics.

The formulation of a closed equation that consolidates these ideas remains a subject for future investigation.

% ----------------------------------------
% Subsection 3.3: Experimental Verification Prospects
% ----------------------------------------
\subsection{Experimental Verification Prospects}

% ----------------------------------------
% Subsubsection 3.3.1: Phase-Dependent Inertial Mass Measurements
% ----------------------------------------
\subsubsection{Phase-Dependent Inertial Mass Measurements}
Should future technology allow for the generation of pulses shorter than the oscillation period of the electron’s internal Zitterbewegung, we could directly test the proposed phase---dependent inertial mass---a fundamentally novel prediction with far-reaching implications for quantum control. By precisely targeting moments when potential energy is completely localized in either Kernel $A$ or $B$, and by appropriately capturing the timing of subsequent radiation transfer, we could potentially reduce the energy required to move electrons.

This prediction arises because the destination point $B$ is determined by the principle of least action. We theorize that by applying force at the precise moment when this destination is being determined, the electron could be moved with significantly less energy input. The technical challenge lies in generating sufficiently short pulses corresponding to the calculated frequency of approximately $5.0 \times 10^{18}$ Hz~\cite{hanamura2023}, and synchronizing them with the electron's internal phase. While currently beyond our experimental capabilities, advances in ultrafast physics suggest such verification experiments may become feasible as technology continues to develop.

If pulses of equal energy produce varying electron behaviors depending on their timing, this would provide further evidence for internal structure and temporal phase-dependent inertial mass. This principle might also enable measurement of electron phase and spin states.

% ----------------------------------------
% Subsubsection 3.3.2: Spin State Manipulation
% ----------------------------------------
\subsubsection{Spin State Manipulation}
Since our model proposes that electron spin alternates between up and down states according to temporal phase, controlling the internal temporal phase could potentially allow deliberate manipulation of electron spin states—opening new possibilities for quantum control at the subatomic level.

If such spin state manipulation were realized through phase control, it could provide a fundamentally new mechanism for qubit initialization and manipulation in quantum computing. Unlike conventional spin control methods that rely on external magnetic fields or spin-orbit coupling, this approach would enable purely internal modulation of spin states by adjusting the electron’s internal energy phase. This would allow for potentially faster, lower-energy, and more localized spin transitions, offering a promising avenue for scalable quantum technologies. Moreover, the ability to directly govern spin via temporal phase could lead to new forms of spin coherence preservation, since such control would bypass many environmental decoherence pathways associated with traditional external-field methods.

This perspective motivates a broader consideration of how spin state control has evolved in quantum physics. Recent progress in quantum state control provides a foundational context for our theoretical proposal. The field has evolved from early spin resonance experiments by Rabi, Ramsey, and others~\cite{rabi1937}, to the development of coherent manipulation techniques for two-level systems~\cite{allen1975}, which established the theoretical underpinnings of quantum control. Experimental milestones, such as precise state engineering in trapped ion systems~\cite{monroe1995}, demonstrated the feasibility of these principles in practice. Within this conceptual framework, our proposed mechanism—controlling electron spin through internal temporal phase modulation—offers a novel approach to deterministic spin state preparation, with potential advantages in speed and energy efficiency over conventional methods relying on thermalization or external driving fields.

% ----------------------------------------
% Subsubsection 3.3.3: Deterministic Control of Quantum Phenomena
% ----------------------------------------
\subsubsection{Deterministic Control of Quantum Phenomena}
Additionally, if technology to measure and control the electron's internal energy transfer could be developed, it might fundamentally transform our understanding of quantum phenomena such as the double-slit experiment. By controlling the distribution of TPE between kernels $A$ and $B$, it might become possible to deliberately direct electrons toward specific regions of the detection screen, steering them left or right as desired. Such control would potentially move quantum mechanics beyond probabilistic descriptions toward deterministic manipulation of quantum behavior. This would represent a paradigm shift in our understanding of quantum phenomena, resolving long-standing philosophical debates about wave-particle duality.

As demonstrated in our previous work~\cite{hanamura2019}, when the electron's phase is at integer multiples of $\pi /2$ (where either $\cos^4\left(\frac{\omega t}{2}\right)=1$ or $\sin^4\left(\frac{\omega t}{2}\right) =1$), one kernel contains all the energy, producing no interference pattern even with both slits open. By selectively controlling this phase, we could switch between wave-like and particle-like behaviors on demand. Such control would potentially move quantum mechanics beyond probabilistic descriptions toward deterministic manipulation of quantum behavior, resolving long-standing philosophical debates about wave-particle duality and addressing Einstein's fundamental objection that ``God does not play dice with the universe.''


% ========================================
% Section 4: Conclusion
% ========================================
\section{Conclusion}

We propose that the origin of electron mass lies in the internal Zitterbewegung (trembling motion) intrinsic to the electron. The oscillation of thermal potential energy between two internal kernels generates spin angular momentum, effectively modeling the electron as possessing gyroscopic properties with a well-defined internal structure. This gyroscopic behavior may constitute the fundamental basis of what we perceive as inertial mass.

This study establishes a conceptual connection between the classical gyroscopic effect and the spin behavior of electrons in relativistic quantum mechanics. By modeling electron spin as a harmonic oscillator subject to sinusoidal acceleration and incorporating geodetic precession, we demonstrate that a spinning electron exhibits resistance to changes in orientation, analogous to the behavior of a macroscopic gyroscope. The combined effects of Thomas precession, Lorentz contraction, and geodetic correction provide a unified framework for understanding spin-induced inertial properties.

Our analysis incorporates both special and general relativistic effects, reaffirming previously published predictions—from a purely special relativistic calculation ($v_{e,\text{SR}} = 0.040472c$) to a refined estimate including general relativistic corrections ($v_{e,\text{SR+GR}} = 0.040374c$). While these values were derived in earlier studies, their inclusion here provides a quantitative reference for the internal dynamics under discussion. The subtle but measurable difference between the two predictions highlights the role of spacetime curvature even at quantum scales.

The model challenges the conventional field-theoretic view and opens a pathway toward deterministic quantum theories, echoing Einstein’s vision of a complete quantum mechanics~\cite{einstein1935}. Future experiments, such as high-precision electron scattering or spin-dependent inertial effect measurements, may be able to test the predicted Zitterbewegung velocity of $0.040374c$. Rather than rejecting the Higgs mechanism~\cite{higgs1964}, this work explores whether an internal geometric structure at the quantum level could contribute to the origin of inertial mass. To further develop this perspective, a quantitative framework such as the derivation of the energy-momentum tensor will be necessary.

The proposed connection between spin and mass implies that quantum-scale dynamics may underlie macroscopic classical phenomena. Our model provides a theoretical framework that unifies quantum mechanics, special relativity (through Lorentz contraction), and general relativity (through geodetic precession), offering a geometric basis for the emergence of mass from more fundamental processes.

In summary, we propose a geometrically grounded interpretation of inertial mass:  
it arises from the internal resistance to reorienting the energy transfer axis embedded within the electron.  
This resistance, governed by the Zitterbewegung dynamics between dual energy kernels, encapsulates the interplay of mass, spin, and inertia as emergent properties of an intrinsic and self-consistent spacetime structure.


% ========================================
% Bibliography
% ========================================
\begin{thebibliography}{99}

\bibitem{hanamura2023} 
S. Hanamura, Redefining Electron Spin and Anomalous Magnetic Moment Through Harmonic Oscillation and Lorentz Contraction, Zenodo (2023), \href{https://doi.org/10.5281/zenodo.17764997}{doi:10.5281/zenodo.17764997}. [Originally: \href{https://vixra.org/abs/2309.0047}{viXra:2309.0047}]

\bibitem{bohm1952} Bohm, D., 
\textit{A Suggested Interpretation of the Quantum Theory in Terms of "Hidden" Variables}, 
Phys. Rev. \textbf{85}, 166 (1952). DOI: \href{https://doi.org/10.1103/PhysRev.85.166}{10.1103/PhysRev.85.166}

\bibitem{schrodinger1930}
Schr\"{o}dinger, E.,
\textit{Über die kräftfreie Bewegung in der relativistischen Quantenmechanik}, 
Sitzungsber. Preuss. Akad. Wiss. Phys.-Math. Kl. \textbf{24}, 418 (1930). 
\href{https://archive.org/details/uberderkrarelativistischenquantenmechanikschrodinger19300717english}{archive.org  (Open Access)}.

\bibitem{hanamura2018} 
S. Hanamura, A Model of an Electron Including Two Perfect Black Bodies, Zenodo (2018), \href{https://doi.org/10.5281/zenodo.16759284}{doi:10.5281/zenodo.16759284}. [Originally: \href{https://vixra.org/abs/1811.0312}{viXra:1811.0312}]

\bibitem{hanamura2024} 
S. Hanamura, Quantum Oscillations in the 0-Sphere Model: Bridging Zitterbewegung, Geodesic Paths, and Proper Time Through Radiative Energy Transfer, Zenodo (2024), \href{https://doi.org/10.5281/zenodo.17765017}{doi:10.5281/zenodo.17765017}. [Originally: \href{https://vixra.org/abs/2411.0117}{viXra:2411.0117}]

\bibitem{hanamura2020} 
S. Hanamura, Correspondence Between a 0-Sphere and the Electron Model, Zenodo (2020), \href{https://doi.org/10.5281/zenodo.17759908}{doi:10.5281/zenodo.17759908}. [Originally: \href{https://vixra.org/abs/2001.0610}{viXra:2001.0610}]

\bibitem{thomas1926} Thomas, L. H., 
\textit{The Motion of the Spinning Electron}, 
Nature \textbf{117}, 514 (1926). DOI: \href{https://doi.org/10.1038/117514a0}{10.1038/117514a0}

\bibitem{thomas1927} Thomas, L. H., 
\textit{The Kinematics of an Electron with an Axis}, 
Philosophical Magazine \textbf{3}, 1 (1927). DOI: \href{https://doi.org/10.1080/14786440108564170}{10.1080/14786440108564170}

\bibitem{weinberg1995} Weinberg, S., 
\textit{The Quantum Theory of Fields}, Cambridge University Press (1995).

% The second paper, written the following March, concerns the static theory of gravitational fields.
\bibitem{einstein1912} Einstein, A.,
\textit{Lichtgeschwindigkeit und Statik des Gravitationsfeldes} [The Speed of Light and the Static Gravitational Field],
Annalen der Physik \textbf{38}, 355-369 (1912). \href{http://myweb.rz.uni-augsburg.de/~eckern/adp/history/einstein-papers/1912_38_355-369.pdf}{Full Text (Open Access)}

\bibitem{fan2023} Fan, X., Myers, T. G., Sukra, B. A. D., Gabrielse, G., 
\textit{Measurement of the Electron Magnetic Moment}, 
Phys. Rev. Lett. \textbf{130}, 071801 (2023). \href{https://arxiv.org/abs/2209.13084}{arXiv:2209.13084}

\bibitem{hartle2021} Hartle, J. B., 
\textit{Gravity - An Introduction to Einstein's General Relativity}, p.302, Cambridge University Press (2021).

\bibitem{hanamura2025} 
S. Hanamura, Bridging Quantum Mechanics and General Relativity: A First-Principles Approach to Anomalous Magnetic Moments and Geodetic Precession, Zenodo (2025), \href{https://doi.org/10.5281/zenodo.17765097}{doi:10.5281/zenodo.17765097}. [Originally: \href{https://vixra.org/abs/2501.0130}{viXra:2501.0130}]

\bibitem{penrose1971} Penrose, R., 
\textit{Angular Momentum: An Approach to Combinatorial Space-Time}, 
in \textit{Quantum Theory and Beyond}, ed. Bastin, T., Cambridge University Press (1971).  \href{https://math.ucr.edu/home/baez/penrose/Penrose-AngularMomentum.pdf}{math.ucr.edu (Open Access)}

\bibitem{mach1919} Mach, E., 
\textit{The Science of Mechanics: A Critical and Historical Account of Its Development}, 
Open Court Publishing (1919). \href{https://archive.org/details/scienceofmechani00machuoft}{archive.org  (Open Access)}

\bibitem{barbour2001} Barbour, J., 
\textit{The Discovery of Dynamics: A Study from a Machian Point of View of the Discovery and the Structure of Dynamical Theories}, 
Oxford University Press (2001). ISBN: 9780195132021

\bibitem{rabi1937} Rabi, I. I., 
\textit{Space Quantization in a Gyrating Magnetic Field},
Physical Review \textbf{51}, 652 (1937).
DOI: \href{https://doi.org/10.1103/PhysRev.51.652}{10.1103/PhysRev.51.652}

\bibitem{allen1975} Allen, L., Eberly, J. H.,
\textit{Optical Resonance and Two-Level Atoms},
Dover Publications (1975).
ISBN: 9780486655338.
\href{https://store.doverpublications.com/0486655334.html}{Publisher's Link}

\bibitem{monroe1995} Monroe, C., Meekhof, D. M., King, B. E., Itano, W. M., Wineland, D. J., 
\textit{Demonstration of a Fundamental Quantum Logic Gate},
Physical Review Letters \textbf{75}, 4714 (1995).
DOI: \href{https://doi.org/10.1103/PhysRevLett.75.4714}{10.1103/PhysRevLett.75.4714}

\bibitem{hanamura2019} 
S. Hanamura, Perfect Contrast Cannot be Obtained in the Electron Double-Slit Experiment, Zenodo (2019), \href{https://doi.org/10.5281/zenodo.17759699}{doi:10.5281/zenodo.17759699}. [Originally: \href{https://vixra.org/abs/1906.0275}{viXra:1906.0275}]

\bibitem{einstein1935} Einstein, A., Podolsky, B., Rosen, N., 
\textit{Can Quantum-Mechanical Description of Physical Reality Be Considered Complete?}, 
Phys. Rev. \textbf{47}, 777 (1935). DOI: \href{https://doi.org/10.1103/PhysRev.47.777}{10.1103/PhysRev.47.777}

\bibitem{higgs1964} Higgs, P. W., 
\textit{Broken Symmetries and the Masses of Gauge Bosons}, 
Phys. Rev. Lett. \textbf{13}, 508 (1964). DOI: \href{https://doi.org/10.1103/PhysRevLett.13.508}{10.1103/PhysRevLett.13.508}

\end{thebibliography}

% = = * * = =  * * = = * * = = * * = = * * = = * * = = * * = = * * 
% \section{\label{sec:Appendix}Appendix}

% \section*{Acknowledgments}

\end{document}